\section{PixArt-$\alpha$：让 DiT 理解自然语言}

Text-to-image 把短类别条件升级为长文本 token。架构必须同时完成图像 patch 内部建模与文本--图像对齐。PixArt-$\alpha$ 的路线是保留 DiT self-attention，用 cross-attention 读取文本，并继续用 timestep 调制 AdaLN。

\subsection{文本条件需要哪些新增模块}

把 prompt 编成单个向量会丢失词级关系，因此主流模型保留文本 token sequence。读下面两页时，应分别追踪：文本在哪编码、在哪与图像交互、时间步又如何进入每层。本节从 Original DiT 的低维类别条件过渡到自然语言条件，问题不只是“多接一个 encoder”，而是怎样保留词级结构、控制接口方向并核算完整 pipeline 成本。

\lecturefigure{slide-35.jpg}{从 class-conditional DiT 到 text-to-image：文本编码、cross-attention 与时间调制}{35}

三个问题构成设计清单：用什么 text encoder；图文通过 cross-attention、拼接还是 joint attention 交互；timestep 是否继续使用 AdaLN。它们可独立选择，因此“text DiT”不是单一架构。

\lecturefigure{slide-36.jpg}{PixArt-$\alpha$ block：self-attention、cross-attention、MLP 与 timestep 条件}{36}

Noisy image tokens 先做 self-attention，再以 image tokens 为 query、text tokens 为 key/value 做 cross-attention。时间步条件通过 AdaLN 调制 image stream。这样文本保持序列结构，图像 token 数与文本 token 数可不同。

\lecturefigure{slide-37.jpg}{Flan-T5-XXL 文本编码器提供长 prompt 与丰富语义理解}{37}

大 text encoder 提升语义能力，却也带来显存与延迟。若 encoder 冻结，生成训练不更新语言表示；若联合训练，成本更高且可能破坏语言能力。比较模型大小时必须说明是否包含 text encoder 参数。

\begin{knowledgebox}{Cross-attention 的方向}
图像 token 作为 query，文本 token 作为 key/value，意味着每个空间位置主动读取相关词义。输出长度仍等于图像 token 数；文本不会被这一层直接更新。MMDiT 会让两种模态共同演化。
\end{knowledgebox}

\subsection{共享 AdaLN 与紧凑模型结果}

PixArt 保留 timestep modulation，但不必为每个 block 都维护独立参数。共享调制器降低计算与参数，同时把每层差异交给 block 自身表示。本节进一步追问这种紧凑性来自哪里：共享的是条件映射而非全部 Transformer 权重，因此必须分别检查参数量、训练计算和最终质量，不能把“共享”自动等同于端到端更快。

\lecturefigure{slide-38.jpg}{PixArt 的 AdaLN：由 timestep 生成 block 调制参数}{38}

代码页显示 time embedding 经 MLP 后产生 scale/shift。与 Original DiT 相同，AdaLN 负责时间条件；文本走 cross-attention。把两类条件拆开，避免把长文本压缩成一个调制向量。

\lecturefigure{slide-39.jpg}{共享 AdaLN 参数并为不同 block 加轻量偏置}{39}

若所有 block 共享同一调制器，可把每层额外参数从 $O(LD^2)$ 降到更接近 $O(D^2+LD)$，其中 $L$ 是层数。共享节省约束了表达能力，但实验显示可在较小质量损失下显著降成本。

\lecturefigure{slide-40.jpg}{PixArt-$\alpha$：紧凑参数量下仍取得有竞争力的文本到图像结果}{40}

表格与柱状图要同时读参数、训练成本与指标。模型更小不等于训练更简单：数据 caption 质量、分辨率 curriculum、VAE、optimizer 和 dropout 都可能贡献结果。讲者明确把 training recipe 留作另一个主题。

\teachervoice{00:32:56--00:40:45，老师把 PixArt 的增量拆成两条：用大文本编码器与 cross-attention 处理自然语言；继续用 timestep AdaLN，并通过共享参数节省计算。架构简洁不代表训练“轻松”。}

\begin{warningbox}{参数量口径必须一致}
报告 denoiser 参数时可能不含冻结 text encoder 与 VAE；端到端部署却必须加载它们。比较“模型更小”时，应同时列 denoiser、condition encoders、峰值显存、训练 FLOPs 和采样延迟。
\end{warningbox}

\subsection{本章小结}

PixArt-$\alpha$ 通过文本 encoder 与 cross-attention 把 DiT 扩展到自然语言，同时保留 timestep AdaLN 并共享调制参数。它展示了紧凑 backbone 的潜力，也提醒我们必须把架构、训练 recipe 与完整 pipeline 成本分开。

